\section{Deletion measures and local transport}
\label{sec:local-palm}

We now prove the adjacent-cell comparison~\eqref{eq:local-fourth} for
the square roots of the deletion densities.  We represent each deletion
measure by an augmented path sample with a distinguished flagged label,
then construct one measure that is close to the two adjacent deletion
measures in fourth moment.  The construction follows the local deletion argument of
\cite[Section~6]{D}; every conditional law below uses the exact
soft-interaction checks of Lemma~\ref{lem:conditional-path-law}.

\subsection{Augmented deletion measures and count errors}

For a cell $v$ of $\mathcal C_d$ or $\mathcal C_R$, define
\[
 T_v=f_0D|v|,\qquad h_{vi}=\1_{\{X_i\in v\}},\qquad
 A_v=\frac1{T_v}\sum_{i\text{ flagged}}h_{vi}.
\]
Here $X_i=\omega_i(0)$, and the expectation $\E$ includes the ground
path sample, its independent flags, and the auxiliary variables in
\eqref{eq:soft-check}.  Define a finite measure on this augmented
sample together with a distinguished flagged label by
\[
 \int H\,\dd P_{\mathrm{aug},v}
   =\E\sum_{i\text{ flagged}}\frac{h_{vi}}{T_v}H(\omega,\cD,i).
\]
Delete the distinguished midpoint and order the other $N-1$ midpoints
uniformly at random, independently of all existing variables.  The
resulting projected measure is $P_v$ from~\eqref{eq:palm-density}.
Indeed the midpoint law is $\Phi^2\dd X$, a specified flag has
probability $f_0$ independently of that law, and symmetry makes all
$N$ tag terms identical after projection.  Their total coefficient
is $Nf_0/T_v=V/|v|$.  The mass of both the augmented measure and its
projection is $\E A_v=1$.  We keep these reference masses explicit
when restricting tags and comparing the finite measures below.

We use normalized averages over cells and over adjacent ordered cell
pairs.  The periodic cubic adjacency graph is regular, so either cell
marginal of its full edge average is a uniform cell average.

\begin{lemma}[Flagged counts and omitted tags]\label{lem:local-counts}
The local energy estimates imply
\begin{equation}\label{eq:flag-moments}
 \begin{aligned}
 \E_{v\in\cC_R}\E|A_v-1|^2
   &\le C\bigl(D^{-1-w/6}+D^{-1+x}R^{-3}\bigr),\\
 \E_{v\in\cC_d}\E|A_v-1|^4
   &\le C_j\bigl(D^{-1-u/6}+D^{-2+2x}\bigr).
 \end{aligned}
\end{equation}
Let
\[
 M_v'=\frac1{T_v}\sum_{i\text{ flagged},\ i\notin I}h_{vi}
\]
be the relative count discarded when only tags in $I$ are retained.
Uniformly over cells of either size,
\begin{equation}\label{eq:tag-omission}
                    \E M_v'\le C_{x,j}D^{-1-1/8}.
\end{equation}
In expectations involving a fixed number of cells, each additional
relative-count factor can be bounded by $C_jD^x$, at an error smaller
than any prescribed power of $D^{-1}$.  The same assertion holds for
fixed products of such factors.
\end{lemma}

\begin{proof}
Given the midpoint configuration, the flagged count $N_v^{\rm f}$
is binomial with parameters $n_v$ and $f_0$.  Thus
\[
 \E[(N_v^{\rm f}-f_0n_v)^2\mid X]\le f_0n_v,
 \qquad
 \E[(N_v^{\rm f}-f_0n_v)^4\mid X]
     \le C\bigl(f_0n_v+(f_0n_v)^2\bigr).
\]
Apply these inequalities to
\[
 A_v-1=\frac{N_v^{\rm f}-f_0n_v}{f_0D|v|}
                  +\left(\frac{n_v}{D|v|}-1\right).
\]
The second summand is controlled by Proposition~\ref{prop:local-energy}.
The binomial contribution to the fourth moment is
$C_j(D^{-2+2x}+D^{-3+3x})$, whose second term is smaller for $x<1$.
The required ordinary count moments follow from
\eqref{eq:count-mgf}.  This proves \eqref{eq:flag-moments}.

For the tail assertion, restrict first to the event that every unit
cell meeting the indicated cells has count at most $CD$, for a
sufficiently large fixed $C$.  Then $A_v\le C_j/f_0=C_jD^x$.
The complement has probability bounded by a polynomial in $D$ times
$e^{-cD}$, even for an $R$-cell, by \eqref{eq:count-tails} and a union
bound over $O(R^3)$ unit cells.  Equation~\eqref{eq:count-mgf} bounds
every fixed moment of the counts there by a polynomial in $D$.
Cauchy--Schwarz therefore makes any fixed count product on this
complement smaller than every prescribed power of $D^{-1}$.

For \eqref{eq:tag-omission}, independence of flagging cancels $f_0$
from the normalization.  The motion failures, including a violation
of \eqref{eq:motion-cutoff} or an ordinary displacement exceeding
$b_0$, have a power saving stronger than
\eqref{eq:tag-omission}, by \eqref{eq:visitor-exponential}.
On their complement a possible tag
in $v$ has both endpoints in a fixed enlargement of $v$, and its tube
condition holds.  At either endpoint time, the expected number of
labels in a unit neighborhood having a second endpoint within $s_1$
is at most $C_xD^2s_1^3$, by \eqref{eq:pair-bound}.  Division by $D$
gives $C_xDs_1^3=C_xD^{-7/5}$.  For an $R$-cell the number of unit
neighborhoods is canceled by $|v|$ in the normalization.  These bounds
are stronger than \eqref{eq:tag-omission}.
\end{proof}

The comparisons will also discard failures of good data near their
endpoint cells.  For adjacent $d$-cells $v,w'$, let
$\mathcal G_{v,w'}$ mean that every site in a fixed neighborhood of
radius $C b_0$ of the two cells is good.  Choose $C$ large enough to
contain every possible tag mean and the coordinate routes used below.
This event depends only on $\cD$.  For two $R$-cells the corresponding
neighborhoods have radius $C_xR$.  By \eqref{eq:good-prob}, their
failure costs, in either augmented reference, at most
\begin{equation}\label{eq:good-omission}
 C_{x,j}D^{-1-1/8+x}\quad\hbox{at scale }d,
 \qquad C_xR^3D^{-1-1/8+x}\quad\hbox{at scale }R.
\end{equation}
To see this, use the union bound for the number of tested sites, and
bound the reference weight by $A_v\le C_jD^x$ off the negligible count
tail.  Multiplication of any omission by a fixed product of relative
counts is handled by the last assertion of Lemma~\ref{lem:local-counts}.
Also $M_v'\le A_v$, so on the same count event
$(M_v')^q\le(C_jD^x)^{q-1}M_v'$.

\subsection{Two copies with the same remaining midpoints}

Fix adjacent $d$-cells $v,w'$ and retain $\mathcal G_{v,w'}$.
Choose distinct tags $i,i_m\in I$ with weight
\begin{equation}\label{eq:tag-pair-weight}
                         \frac{h_{vi}h_{w'i_m}}{T_vT_{w'}}.
\end{equation}
The site of a tag is the unit cell containing its near endpoint mean
$m_i$.  The tube condition places this site in a fixed enlargement of
the indicated cell whenever its tag indicator is nonzero.  Join the
two tag sites by a simple coordinate route of nearest-neighbor sites.
At each interior vertex choose an interior label uniformly from the
available supply, excluding the two tags.  Choices at different sites
are independent given $\cD$ and the tags.  This gives a list
\begin{equation}\label{eq:local-list}
                 i_0=i,i_1,\ldots,i_m,
 \qquad m\le C_x=O(x^{-1/2}).
\end{equation}
If the two tag sites coincide, take just the two tags.  The labels are
distinct, consecutive means are at an absolute bounded distance, and
each interior choice has point probabilities at most $C_x/(Df_0)$.
The route and these probability kernels depend only on the data and
the tag identities; the weight \eqref{eq:tag-pair-weight} does not
alter their definition.

Construct a left and a right path configuration.  Keep all unlisted
paths in both copies, the original path of $i_0$ on the left, and the
original path of $i_m$ on the right.  For $k=1,\ldots,m$, generate
the left path of $i_k$ together with the right path of $i_{k-1}$.
Choose a common midpoint uniformly in a unit cell $C_k$ at an absolute
bounded distance from both means.  Fix $C_k$ as a measurable function
of the endpoints before drawing any paths.  Given the midpoint,
propose the two paths independently as pinned free bridges.  Condition
this entire joint proposal, including its random midpoint, on the
following tests: each path satisfies its
individual cutoffs and every check \eqref{eq:soft-check} against
retained paths and paths already inserted in its own copy.
Proposition~\ref{prop:insertion}(2) gives success probability at least
$1/2$, so every stage is a probability kernel.  The accepted midpoint
need not be uniformly distributed.

Every interaction in a finished copy is checked when the later of its
two paths is inserted, or was already checked between retained paths.
Both copies therefore belong to the same conditional data fiber.
Their midpoint identities are
\[
                X_{i_k}^{\rm left}=X_{i_{k-1}}^{\rm right},
                         \qquad 1\le k\le m.
\]
Deleting $i_0$ on the left and $i_m$ on the right leaves the same
midpoint multiset.  Give the left remaining labels a uniform random
order and give the right labels the order induced by these pairings
and the unchanged labels.  Each order is marginally uniform.
Weighting the original sample by \eqref{eq:tag-pair-weight}, and
averaging tags, choices, and these kernels, defines a single finite
measure $\nu_{v,w'}$ on the remaining coordinates.  It has a lift to
either augmented reference: observe the left copy with tag $i_0$, or
the right copy with tag $i_m$.
Figure~\ref{fig:local-transport} displays the midpoint pairing.

\begin{figure}[tb]
\centering
\begin{tikzpicture}[x=1.65cm,y=1.3cm,>=Latex,
 particle/.style={circle,draw,minimum size=8mm,inner sep=1pt},
 made/.style={particle,fill=blue!12,draw=blue!65!black},
 tag/.style={particle,fill=gray!12,draw=gray!65!black}]
\node[anchor=east] at (-.55,1) {left copy};
\node[anchor=east] at (-.55,0) {right copy};
\foreach \k in {0,...,3}{
 \ifnum\k=0
   \node[tag] (L\k) at (\k,1) {$i_\k$};
 \else
   \node[made] (L\k) at (\k,1) {$i_\k$};
 \fi
 \ifnum\k=3
   \node[tag] (R\k) at (\k,0) {$i_\k$};
 \else
   \node[made] (R\k) at (\k,0) {$i_\k$};
 \fi
}
\foreach \k/\j in {1/0,2/1,3/2}{
 \draw[blue!65!black,thick] (L\k) -- node[fill=white,inner sep=1pt,font=\small] {$y_\k$} (R\j);
}
\draw[red!65!black,thick] (L0.south west)--(L0.north east);
\draw[red!65!black,thick] (R3.south west)--(R3.north east);
\node[font=\small,anchor=south] at (0,1.4) {delete};
\node[font=\small,anchor=north] at (3,-.4) {delete};
\end{tikzpicture}
\caption{Midpoint transport along a three-step label list. The left path
of $i_k$ and the right path of $i_{k-1}$ are generated with common
midpoint $y_k$. Their full paths need not agree. After deleting $i_0$
on the left and $i_3$ on the right, the displayed midpoint multisets
coincide. Unlisted paths are retained in both copies.}
\label{fig:local-transport}
\end{figure}


We calculate a lift density before estimating it.  Treat the left as
the observed copy, fix $\cD$, the tags and list, and condition on the
paths $\eta$ outside
$S=\{i_1,\ldots,i_m\}$.  Put $B_S=\prod_{l\in S}B_l$.  Let
$c_S(o;\eta)$ be the product of all individual cutoff indicators on
$S$ and all checks involving at least one label in $S$.  By
Lemma~\ref{lem:conditional-path-law}, the original conditional law is
\begin{equation}\label{eq:local-reference}
 \lambda_S(\dd o\mid\eta)
       =\frac{c_S(o;\eta)}{Z_S(\eta)}B_S(\dd o),
 \qquad Z_S(\eta)=\int c_S(o;\eta)B_S(\dd o)>0.
\end{equation}
Checks wholly outside $S$ are fixed and do not enter this formula.

The generation kernel on the changed left paths has a density
$k_S(o,\eta,z)$ with respect to $B_S(\dd z)$.  To verify this
precisely, write $B_l^y$ for the free path law of label $l$ conditional
on midpoint $y$, and $g_l$ for its Gaussian midpoint density.  A pair
with observed label $l$ and auxiliary label $l'$ has proposal measure
\begin{equation}\label{eq:local-proposal}
 \frac{\1_C(z_l(0))}{g_l(z_l(0))}
                  B_l(\dd z_l)B_{l'}^{z_l(0)}(\dd z'_{l'}).
\end{equation}
Multiplication by its success indicators and division by its success
probability describes the actual insertion kernel.  Sequential
integration of the auxiliary paths gives $k_S$.  This is a density
for the observed paths only: the paired paths have a common midpoint,
and their joint law has no density against two independent bridges.

\begin{lemma}[The centered conditional likelihood]
\label{lem:local-likelihood}
Include the second tag indicator but omit its divisor $T_{w'}$.
On admissible output paths the left likelihood is
\begin{equation}\label{eq:local-likelihood}
 L_S(z;\eta)=\int c_S(o;\eta)h_{w'i_m}(o)
                      k_S(o,\eta,z)B_S(\dd o).
\end{equation}
It obeys
\begin{equation}\label{eq:local-centering}
 F_S=L_S-h_{w'i_m},\qquad
                 \int F_S\,\dd\lambda_S(\cdot\mid\eta)=0,
\end{equation}
and
\begin{equation}\label{eq:local-likelihood-bound}
                     0\le L_S\le C_xD^{C\sqrt x}.
\end{equation}
These assertions remain valid after deleting specified outside checks
consistently from the reference, the original-path integral, and all
insertion tests and their normalizing probabilities.
\end{lemma}

\begin{proof}
Before taking the likelihood, the generated density against $B_S$
is the right side of \eqref{eq:local-likelihood} divided by $Z_S$.
The output kernel is supported on $c_S(z;\eta)=1$, where the reference
density against $B_S$ is $1/Z_S$.  Thus the normalizer cancels.
Since each generation kernel has mass one, its support also gives
\[
 \int c_S(z;\eta)L_S(z;\eta)B_S(\dd z)
       =\int c_S(o;\eta)h_{w'i_m}(o)B_S(\dd o).
\]
Division by $Z_S$ proves \eqref{eq:local-centering}.  The subtracted
indicator there is evaluated on the reference output, whereas the
indicator inside \eqref{eq:local-likelihood} is evaluated on the
integrated original path.

On $C_k$, the factor $g_{i_k}^{-1}$ in
\eqref{eq:local-proposal} is at most $CD^{Cx}$, since its distance
from the mean is absolutely bounded and $t^{-1}=x\log D$.
Every success denominator costs at most two.  Dropping unused checks,
and integrating the auxiliary conditional bridges, bounds the observed
kernel density by $(CD^{Cx})^m2^m$.  The original-path integral has
mass at most one.  Since $m=O(x^{-1/2})$, this proves
\eqref{eq:local-likelihood-bound}.

Deleting checks preserves output support in the corresponding
admissible set and only improves the insertion lower bounds.  The
same argument applies on that modified conditional fiber.  For later
integrals we also define these kernels at outside values satisfying
all individual cutoffs, even if a check wholly outside $S$ fails:
that outside-only check is ignored.  The insertion bounds hold at
these histories, and the formulas represent actual conditional
likelihoods wherever the outside reference constraints hold.
\end{proof}

\subsection{Cancellation in a fourth moment}

The likelihood for one chosen list need not be close to one.  Averaging
many possible tags and interior labels produces the needed small
error.  Conditional centering cancels terms involving an isolated
list; the terms that survive have either repeated labels or at least
two rare interactions between lists.  We make both statements precise.

For fixed tags and data, put a bar over a quantity to denote averaging
over the route and interior label choices.  We claim that, for
sufficiently small fixed $x$,
\begin{equation}\label{eq:local-centered-fourth}
 \E\left[\1_{\mathcal G_{v,w'}}
       \sum_{i\in I}\frac{h_{vi}}{T_v}
       \left|\sum_{i_m\ne i}\frac{\overline{F_S}}{T_{w'}}\right|^4
       \right]\le C_{x,j}D^{-1-1/20}.
\end{equation}
Each tag sum is restricted to the fixed enlargement of its cell and
to labels in $I$.  We first prove the conditional interaction
estimate needed to expand this fourth moment.

Fix four lists with the same first tag $i$, and let $S_a$ be their
changed sets, $1\le a\le4$.  Condition on $\cD$ and on all paths
outside $U=\bigcup_aS_a$.  Join two lists if their changed sets
intersect; the connected components are called blocks.  A singleton
block contains one list, not necessarily one label.  Distinct blocks
have disjoint label sets.  Introduce a switch
$\sigma_e\in\{0,1\}$ for each unordered pair $e$ of blocks.
Turning it off deletes all checks between those blocks from the
following three places:
\begin{enumerate}
\item the union reference on $U$;
\item the original-path indicators in each likelihood integral;
\item the insertion indicators and their success probabilities.
\end{enumerate}
In the last two places, paths outside a given changed set are the
reference paths.  All individual cutoffs are retained.  The unchanged
left path of the common tag $i$ is outside every block, so checks
against it are never switched.  A newly generated right path with
label $i$, however, belongs to its generating block; its checks
against other blocks are switched.  It is never checked against the
left copy of the same label.  These rules also specify the switches
when two lists overlap.

Let $c_U^\sigma$ denote the switched union indicator and
$F_a^\sigma$ the switched centered likelihood.  Define the
unnormalized integral
\begin{equation}\label{eq:local-switch-integral}
 G(\sigma)=\int c_U^\sigma(z)
                   \prod_{a=1}^4F_a^\sigma(z)B_U(\dd z).
\end{equation}
The conditional fourth-moment summand before switching is
$G(\boldsymbol1)/Z_U$, where $Z_U=\int c_U^{\boldsymbol1}\dd B_U$.
We hold this denominator fixed throughout the switch expansion.
Successive single-path insertions give $Z_U^{-1}\le C_x$, because
$U$ has at most $4C_x$ labels.

For a block-pair edge $e$, let $\mathcal Q_e$ be the finite list of
distinct-label pairs whose checks are controlled by that switch.
It includes checks using generated auxiliary paths and original
paths, with bounded multiplicity.  All its labels and endpoints are
fixed by the data and the four lists.

\begin{lemma}[One and two switched interactions]
\label{lem:local-switches}
Suppose the four changed sets are either pairwise disjoint, or have
exactly one independent label identification.  Write $\Delta_e$ for
the finite difference from $\sigma_e=0$ to $\sigma_e=1$.  Uniformly
in all other switch settings,
\begin{equation}\label{eq:local-one-switch}
 \frac{|\Delta_eG|}{Z_U}
       \le C_xD^{C\sqrt x}\sum_{(l,k)\in\mathcal Q_e}\beta_{lk}.
\end{equation}
If the changed sets are pairwise disjoint and $e\ne f$, then
\begin{equation}\label{eq:local-two-switches}
 \frac{|\Delta_e\Delta_fG|}{Z_U}
   \le C_xD^{C\sqrt x}
       \sum_{\substack{(l,k)\in\mathcal Q_e\\
                       (l',k')\in\mathcal Q_f}}
                                 \beta_{lk}\beta_{l'k'}.
\end{equation}
The sums have at most $C_x$ terms.  Moreover
$\beta_{lk}\le C_xD^{-19/100}$ for sufficiently large $D$.
\end{lemma}

\begin{proof}
Expand each centered factor as $L_S-h_{w'i_m}$, then use
\eqref{eq:local-likelihood} for each likelihood.  Each such term has
reference output bridges on $U$, a separate collection of latent
original bridges for each likelihood, and its generated auxiliary
bridges.  The original integrations remain separate even when their
output coordinates coincide.  All individual cutoff indicators
remain in place when estimating collisions.

First consider singleton blocks.  Combine the observed bridge factor
in \eqref{eq:local-proposal} with its midpoint density and auxiliary
conditional bridge.  The result is exactly the common-uniform-midpoint
pair proposal.  After unused checks and bounded success denominators
are discarded, the variables within a singleton block therefore have
a product of ordinary original bridges and these pair proposals.
If its centered factor contributes the subtracted indicator, it uses
ordinary output bridges instead.  These base measures are independent
between blocks: their cells and endpoints depend only on $\cD$.
The common midpoint inside one pair is preserved throughout.

A difference of a check indicator forces failure of one of the checks
controlled by that edge.  Failure of \eqref{eq:soft-check} requires
a collision within the interaction range $r$.  A union bound reduces
it to one of the label pairs in $\mathcal Q_e$.  Given an allowed
opponent path, integrating a plain bridge or a marginal of the
uniform-midpoint pair proposal, with its individual cutoffs retained,
costs at most $C_x\beta_{lk}$ by \eqref{eq:beta} and the
uniform-midpoint form of \eqref{eq:midpoint-cost}.

Success denominators require the same conclusion.  At fixed history
a success probability can be written using a fresh proposal $Q$ as
\[
 p(\sigma)=\int Q(\dd z)\,c(z)
                         \prod_e A_e(z)^{\sigma_e},
                  \qquad p(\sigma)\ge\tfrac12.
\]
Here $c$ contains the individual cutoffs and unswitched checks, while
$A_e$ is the product of the checks controlled by $e$ at this stage.
We use $A^0=1$ also when $A=0$.  A first difference introduces
$A_e-1$; a mixed difference introduces
$(A_e-1)(A_f-1)$.  They force the indicated one or two collisions.
The identity
\[
                 p^{-1}-q^{-1}=\frac{q-p}{pq}
\]
and its difference in a second variable bound the second reciprocal
difference by a constant times
\[
 |\Delta_e\Delta_fp|
       +\max_\sigma|\Delta_ep(\sigma)|
                       \max_\sigma|\Delta_fp(\sigma)|.
\]
The product of first differences is represented with independent
fresh proposals.  Every such proposal belongs to the block whose
denominator produced it.  The lower bound for $p$ holds at all
histories under consideration, since only individual cutoffs are
required of earlier paths when unused checks are ignored.

The finite-product difference rule now bounds a first difference by
a finite sum of integrals with one forced cross-block collision,
and a second difference by ones with two specified cross-block
collisions.  Unused denominators cost only $C_x$.  With no label
identification, either end of one edge can be integrated as a
singleton.  For two distinct edges, their graph consists of a
two-edge tree or two disjoint edges.  In the tree case fix all
variables at the central vertex and integrate its two leaves in
succession.  Each leaf participates in one forced collision, so these
integrations give the two $\beta$ factors.  No independence between
the two collision events at the central block is required.  For
disjoint edges the two integrations are separate.  Fresh proposals
from denominator differences obey the same argument after assigning
them to their generating blocks.

With one independent identification the block sizes are $2,1,1$.
Every edge has a singleton end, which supplies the collision factor
in \eqref{eq:local-one-switch}.  After integrating that singleton,
the remaining repeated output coordinates in the block of size two
are bounded by the observed density estimate
\eqref{eq:local-likelihood-bound}: bound any surplus factors
$g_l^{-1}\1_C$ by $CD^{Cx}$ and integrate auxiliary conditional
bridges.  Across the four lists this costs at most
$C_xD^{C\sqrt x}$.  The same harmless bound covers all remaining
factors in the disjoint case.  This proves both difference estimates.
The argument allows arbitrary settings of unused switches; additional
differences are bounded by finite sums of these estimates at those
settings.  Finally \eqref{eq:beta} gives
$\beta_{lk}\le b_D(1+2/s_1)\le C_xD^{-19/100}$.
\end{proof}

We can now prove \eqref{eq:local-centered-fourth}.  Expanding the
fourth power gives four independently chosen lists, conditional on
the common first tag and the data.  There are at most $C_xDf_0$
candidates for each tag or interior slot, by the good endpoint
counts and the tube restriction.  A normalized tag counting weight
and every interior-choice probability are at most $C_{x,j}/(Df_0)$.
Split the sums by the tag mean-sites and the route site patterns;
there are only $C_x$ patterns.  We retain the original weights in
each subsum, without conditioning on a pattern and renormalizing.

An equality pattern with $q$ independent identifications has $q$
fewer free label choices.  Dominating all slot weights by their
point bounds and counting the remaining choices gives total weight
at most $C_{x,j}(Df_0)^{-q}$.  A triple equality counts as two
identifications.  These are joint counting estimates; they assert
no point-mass bound conditional on the occurrence of a coincidence.
There are only $C_x$ equality patterns.  Consequently all patterns
with $q\ge2$ contribute, by \eqref{eq:local-likelihood-bound}, at most
\begin{equation}\label{eq:local-many-identifications}
                         C_{x,j}D^{-2+2x+C\sqrt x}.
\end{equation}

For $q=0$ or $1$, use the switches above.  For an edge set $E$, let
$\Delta_E=\prod_{e\in E}\Delta_e$.  The elementary Boolean identity is
\begin{equation}\label{eq:local-boolean}
                   G(\boldsymbol1)=\sum_E\Delta_EG(\boldsymbol0).
\end{equation}
If the graph with edge set $E$ has an isolated singleton block, its
term vanishes.  Indeed all edges from that block are off in every
evaluation comprising the difference.  Its variables separate from
the others, and their integral is
\[
                       \int c_S F_S\,\dd B_S=0
\]
by \eqref{eq:local-centering}.  This is why the common denominator
$Z_U$ was held fixed rather than switched.

When $q=1$ there are three blocks, two of them singletons, so a
nonzero term must contain an edge.  Its identification cost and
\eqref{eq:local-one-switch} give total contribution at most
\begin{equation}\label{eq:local-one-identification}
                    C_{x,j}D^{-1+x-19/100+C\sqrt x}.
\end{equation}
When $q=0$, all four blocks are singletons and a nonzero term has
at least two distinct edges.  Apply \eqref{eq:local-two-switches}
to any two of them, and bound any remaining differences by a finite
sum of evaluations with the other switches fixed.

One of these two collision factors can be averaged over a label
slot.  Every cross-block check has at least one label in a changed
set, whereas the common tag $i$ belongs to none of those sets.
The endpoint bound \eqref{eq:flag-count}, summed over dyadic annuli
in a fixed neighborhood, gives uniformly in the other endpoint $y$
\begin{equation}\label{eq:label-average}
 \frac{C_{x,j}}{Df_0}
       \sum_{k\text{ in a slot range}}
                 \frac1{s_1\vee d_K(y,Y_k)}
 \le C_{x,j}\left(1+\frac{D^{1/100}}{Df_0s_1}\right)
 \le C_{x,j}.
\end{equation}
For example, a dyadic annulus of radius $h\le1$ contributes at most
$C_{x,j}(h^2+D^{1/100}/(Df_0h))$; the innermost ball is bounded at
radius $s_1$, and the unit-scale and larger annuli use the unit-count
bound.  The geometric sums give the displayed expression.  The same
estimate holds with $Z_k$ in place of $Y_k$.  Thus averaging this
factor costs at most $C_{x,j}b_D$.  Bound the other factor by
$C_xD^{-19/100}$.  In these positive bounds we may remove all further
noncoincidence restrictions and enlarge the label sums.  The total
contribution with no identification is therefore at most
\begin{equation}\label{eq:local-no-identification}
                    C_{x,j}D^{-19/100+C\sqrt x}b_D.
\end{equation}
The first-tag factor $h_{vi}$ is fixed in every conditional integral,
so it does not disturb centering.  Its total normalized tag sum on
good data is bounded by $C_{x,j}$.  Equations
\eqref{eq:local-many-identifications}--\eqref{eq:local-no-identification}
prove \eqref{eq:local-centered-fourth} once $x$ is chosen sufficiently
small.  Indeed $b_D=C_xD^{-1}(\log D)^C$, and all three powers then
save more than $1/20$ beyond $D^{-1}$.

\subsection{The local comparison after projection}

Let $\ell_v$ and $\ell_{w'}$ be the two lifted densities of
$\nu_{v,w'}$ relative to their augmented references.  On retained
data with first tag in $I$, the left density is
\[
                  \ell_v=\sum_{i_m\ne i}\frac{\overline{L_S}}{T_{w'}};
\]
set it to zero on the omitted part of the reference.  The centered
subtraction sums to $A_{w'}-M_{w'}'$.  In fact the candidate mean
range contains every tag with nonzero $h_{w'i_m}$, by confinement,
and the excluded first tag has $h_{w'i}=0$ because the cells are
disjoint.  Thus on the retained part
\begin{equation}\label{eq:local-lift-error}
 \ell_v-1=\sum_{i_m\ne i}\frac{\overline{F_S}}{T_{w'}}
                              +(A_{w'}-1)-M_{w'}'.
\end{equation}

\begin{proposition}[Adjacent deletion measures]\label{prop:local-palm}
For sufficiently small fixed $x>0$, depending on $u$, the common
measure constructed above satisfies
\begin{equation}\label{eq:local-lift-fourth}
 \E_{v\sim w'}\left[
       \int|\ell_v-1|^4\,\dd P_{\mathrm{aug},v}
       +\int|\ell_{w'}-1|^4\,\dd P_{\mathrm{aug},w'}\right]
                         \le C_jD^{-1-u/8}.
\end{equation}
Consequently the local comparison~\eqref{eq:local-fourth} holds.
That bound, with a fixed change of constant, also holds when only
adjacent pairs inside the same $R$-cell are retained in the average.
\end{proposition}

\begin{proof}
Apply $|a+b+c|^4\le27(|a|^4+|b|^4+|c|^4)$ to
\eqref{eq:local-lift-error}.  The centered term is controlled by
\eqref{eq:local-centered-fourth}.  Off the negligible count tail,
the count-fluctuation term is at most
\[
 \E[A_v|A_{w'}-1|^4]
     \le C_jD^x\E|A_{w'}-1|^4+O(D^{-A})
\]
for any prescribed fixed $A$.  Averaging adjacent cells and using
\eqref{eq:flag-moments} bounds it by
$C_j(D^{-1-u/6+x}+D^{-2+3x})$.  Likewise,
\[
 \E[A_v(M_{w'}')^4]
       \le C_jD^{4x}\E M_{w'}'+O(D^{-A})
       \le C_{x,j}D^{-1-1/8+4x}+O(D^{-A}).
\]
The part where $\ell_v$ was set to zero contributes its reference
mass, controlled by \eqref{eq:tag-omission} and
\eqref{eq:good-omission}.  Taking $x<u/24$ in addition to the earlier
smallness conditions makes all these terms bounded by
$C_jD^{-1-u/8}$.

For the right lift, regard the right copy as observed, reverse the
list indices, and keep the actual generation order.  Its unchanged
tag is now $i_m$ and its changed set excludes $i_m$.  The density
formula, switched centering, and collision integrations are unchanged:
they allow any fixed insertion order.  This proves
\eqref{eq:local-lift-fourth} for the two lifts of the same measure.

Under a measurable projection, the projected likelihood is the
conditional expectation of the lifted likelihood.  Conditional
Jensen therefore transfers each fourth-moment bound to the density
of $\nu_{v,w'}$ relative to $P_v$ or $P_{w'}$.  Write $q(Y')$ for
its Lebesgue density and put $y=s_v(Y')$, $z=s_{w'}(Y')$.
Absolute continuity with respect to both references makes $q=0$
where either reference density vanishes.  Inserting $\sqrt q$ between
$y$ and $z$, and using $|1-\sqrt a|\le|1-a|$ for $a\ge0$, gives
\[
 \frac{|y-z|^4}{y^2+z^2}
   \le 8\left[y^2\left|1-\frac q{y^2}\right|^4
               +z^2\left|1-\frac q{z^2}\right|^4\right].
\]
Terms with zero prefactor are omitted.  Integrating this inequality
and applying the two projected bounds proves \eqref{eq:local-fourth}.
Restricting the nonnegative edge sum to edges within $R$-cells only
decreases it; the normalizations differ by a bounded factor.
\end{proof}
